\section{Exact evaluation of both complementary-minor Bessel periods}
\label{sec:bessel-period-evaluation-v51}

This chapter proves the whole original \texttt{EvalPStatement} and
\texttt{EvalMStatement}.  It identifies explicit choices of the
standard and dual input vectors and computes their actual periods as
nonzero scalar multiples of the old period.  The character, compact
representation, Haar measures and both free translates remain the
original ones.  The shifted-quotient constructions are still separate.

Fix arbitrary data in those two source statements.  Explicitly,
$n=2r+k+1$, $p=k+r$, $q=r+1$, the original
\texttt{StepGap} holds for $h$, $\lambda$ is Harish--Chandra with
noncompact gaps at least $h$, and
$(\mathcal H,\sigma,D,E,V,T,\mu,\iota)$ is any original complete
unitary derived datum.  It has the prescribed compact irreducible
lowest type, Casimir, actual compact-group stability,
$\mu=\lambda+\rho$, and the exact $\Psi$ coefficient identity.
We impose no irreducibility on the ambient representation.
Let $S=U(\mathfrak{gl}_n)V$, $G=U(J_\epsilon)$,
$G_H=U(H_{pk})$ and $d=n-1=2r+k$.

The original Bessel datum consists of a negative middle vector $z$,
$T_z=-z^*H_{pk}z>0$, nonzero coefficients $c_a$ for $0\le a<r$,
an arbitrary permitted $\pi$ of compact type $u$, and the original
weak Haar measures $\mu_W,\mu_N$.  Fix its Cayley matrix $P$,
any $x_0\in G$, $w\in V$, $\xi\in V_\pi$, and any $g\in G_H$.
For $m\in S$ put
\[
 \begin{split}
 F_m(y)&=\langle w,\sigma(x_0^{-1}y)m\rangle,\\
 \ell(m)&=\int_{H_z\times N}
       F_m(P h'n'gP^{-1})\,
       \overline{\psi_{z,c}(n')}\,\pi(h')^{-1}\xi\,
                    d\mu_W(h')d\mu_N(n').
 \end{split}
 \tag{BSE.1}
\]
This is the literal \texttt{coefL}/\texttt{per} functional.
BST3 and BSP5 prove its absolute convergence and linearity on $S$,
and justify all differentiations used below.  Those results also
apply to every matrix-unit word in $m$.
For the two tensor maps use BST's $Q_f,Q_f^\vee$ with
\[
 x_*= \begin{cases}e_0,&r>0,\\ z,&r=0,\end{cases}
                  \qquad x_1=P x_* .
 \tag{BSE.2}
\]
Their periods are absolutely convergent by BST5 and BSP5.
The arbitrary $g$ above is the point at which each period is evaluated.

\subsection{The actual period and its framed infinitesimal law}

% BEGIN BSE1
\paragraph{BSE1. Exact period equivariance with the original Haar measures.}
\label{claim:BSE1}
For every original Bessel datum, every complex matrix $P$,
every matrix function $F$, every $\xi\in V_\pi$ and every matrix
$x$, the source's period obeys
\[
 \mathcal P_F(h_0m_0x)
       =\psi(m_0)\pi(h_0)\mathcal P_F(x),
       \qquad h_0\in H_z,\quad m_0\in N.
 \tag{BSE.3}
\]
This holds in particular for every function and derivative used in
(BSE.1), and with no scalar change in either Haar measure.

\begin{proof}
By AMD8 the original supported weak Haar measures are ordinary Haar
measures on their closed subgroups.  The group $H_z$ is compact by
BWP3; $N$ is nilpotent and unimodular by the polynomial structure
proved in BWP1 and BNV11.  In exponential coordinates the adjoint
operators of its Lie algebra are nilpotent and have trace zero,
so the modulus of every inner translation is one.  Thus both
measures are also right invariant.
The action of $H_z$ on $N$ preserves the latter Haar measure:
its modulus is a continuous homomorphism from the compact group
$H_z$ into $\mathbb R_{>0}$ and hence equals one.
These arguments concern the given measures themselves, not a
replacement normalization.

The group definitions give $H_zNH_z^{-1}=N$ and
$\psi(h^{-1}nh)=\psi(n)$.  To check the character equality directly,
the first $r-1$ relevant matrix entries are unchanged because
$h$ fixes every outer flag vector, and the last relevant entry
is the coefficient of $nz$, unchanged because $hz=z$.
The same first-level entries add under multiplication in $N$,
so $\psi$ is its actual unitary character.

In the period at $h_0m_0x$ make the measure-preserving substitution
\[
             h''=h'h_0,\qquad
             n''=h_0^{-1}n'h_0m_0.
 \tag{BSE.4}
\]
Then $h'n'h_0m_0=h''n''$,
$\psi(n')^{-1}=\psi(m_0)\psi(n'')^{-1}$, and
$\pi(h')^{-1}=\pi(h_0)\pi(h'')^{-1}$.
Pull the fixed scalar and fixed unitary operator outside the
integral.  This gives (BSE.3).
For the functions used here every interchange is legitimate by
BST3/BST5 and BSP5.  The source states the covariance for arbitrary
matrix functions as well, with its zero convention for undefined
Bochner integrals: the same measurable invertible substitution and
fixed isometry preserve integrability and measurability, so if these
fail both sides are zero.  Thus its unrestricted convention causes
no change to the identity.
\end{proof}
% END BSE1

% BEGIN BSE2
\paragraph{BSE2. The exact compact Capelli polynomial.}
\label{claim:BSE2}
Let $\tau$ be any original $\texttt{IsHCTypeU}(k,u)$ representation.
Its complexified derived matrix units $D_{ij}$ satisfy the
matrix-unit relations and, on all of $V_\pi$,
\[
 \mathcal C_D(t)=p_u(t)I,\qquad
 p_u(t)=\prod_{j=0}^{k-1}
               \left(t+u_j-\frac{k-1}{2}\right).
 \tag{BSE.5}
\]
For every skew-Hermitian $A$, the derivative of
$\tau(\exp(tA))$ is its derived action.  The statement includes $k=0$.

\begin{proof}
The original compact-type predicate is a nonzero finite-dimensional
unitary representation, irreducible of the specified highest weight.
A continuous finite-dimensional Lie-group representation is smooth.
One may see this without an additional classification theorem:
convolution by a smooth approximate identity gives smooth orbit
vectors and an operator arbitrarily close to the identity on the
entire finite-dimensional space; that operator is invertible,
so all its vectors are smooth.  Differentiate the representation
law and complexify $\mathfrak u(k)$ to obtain $D$ and its
matrix-unit relations.

The original compact-type predicate supplies a nonzero torus
vector of weight $\nu_j=u_j-(k-1)/2+j$ and says that every weight
is dominated by $\nu$.  Differentiating torus conjugation shows
that $D_{ab}$ sends that vector to weight $\nu+e_a-e_b$.
If $a<b$, this weight is not dominated by $\nu$: the linear
functional $\sum_j(k-1-j)\mu_j$ strictly increases under adding
$e_a-e_b$ and weakly decreases under subtracting positive roots.
Thus every upper-triangular $D_{ab}$ kills the vector, and it is
the required highest vector.
In the original ordered column determinant
for $\mathcal C_D(t)$, apply the rightmost factor to this vector.
Any nonidentity permutation has a largest column not fixed by it;
the earlier fixed larger columns have already acted diagonally, and
that column's row is strictly smaller than the column.  Its
upper-triangular matrix unit kills the highest vector.  Thus only
the identity permutation survives, giving
$\prod_j(t+\nu_j-j)=p_u(t)$.

The operator $\mathcal C_D(t)$ is central.  This follows directly
from BST6's minor commutator with full identical row and column
lists: if $a\ne b$, each replacement creates a repeated index and
vanishes; if $a=b$, the equal row and column terms cancel.
Hence the scalar identity extends from the highest vector to its
generated module.  That module is the entire irreducible compact
type: its finite-dimensional Lie-stable subspace is invariant
under exponentials, and the connected group $U(k)$ is generated
by them.  Irreducibility completes the argument.
If $k=0$, both the determinant and product in (BSE.5) are the
empty product one, so the assertion remains literal.
\end{proof}
% END BSE2

% BEGIN BSE3
\paragraph{BSE3. The whole original framed period law, with explicit character values.}
\label{claim:BSE3}
Choose a frame $g_0$ furnished by the original
$\texttt{IsHCTypeW}(k,r,z,u,\pi)$.
It is $H_{pk}$-unitary, fixes the outer flag vectors, and sends
$\sqrt{T_z}e_{r+k}$ to $z$.  Put
\[
 Q=g_0^{-1}gP^{-1},\qquad
 E^Q_{ij}=\sum_{a,b}(Q^{-1})_{ai}Q_{jb}E_{ab}.
 \tag{BSE.6}
\]
Let $\tau(A)=\pi(g_0\operatorname{blockW}(A)g_0^{-1})$ on
$U(k)$, and let $D_{ij}$ be its derived units from BSE2.
For the levels
\[
 b(i)=
 \begin{cases}
 i,&i<r,\\
 r,&r\le i\le r+k,\\
 i-k,&i>r+k,
 \end{cases}
 \qquad W=\{r,\ldots,r+k-1\},
 \tag{BSE.7}
\]
there are scalars $\chi_{ij}$ such that, on every $m\in S$,
\[
 \ell(E^Q_{ij}m)=\chi_{ij}\ell(m)\quad(b(i)<b(j)),\qquad
 \ell(E^Q_{r+i,r+j}m)=D_{ij}\ell(m)\quad(i,j<k).
 \tag{BSE.8}
\]
The scalars vanish except on the consecutive edges of
\[
                   0,1,\ldots,r-1,r+k,r+k+1,\ldots,d.
 \tag{BSE.9}
\]
For $r>0$ their product on these edges is exactly
\[
              \kappa=-\frac{T_z}{4^r}
                               \prod_{a=0}^{r-1}|c_a|^2\ne0.
 \tag{BSE.10}
\]
Together with (BSE.5), these are every conclusion of the original
\texttt{FrameLawStatement}.

\begin{proof}
The frame and compact representation are the actual existential
witnesses in the given compact-type predicate.
For a matrix $Z$ write $E(Z)=\sum Z_{ab}E_{ab}$.
The operators $E^Q$ are $E(Q^{-1}e_{ij}Q)$ and preserve $S$;
conjugation of matrices and the original bracket identity show that
they satisfy the same matrix-unit relations.
For a real $H_{pk}$-Lie direction $Y$, the curve through the period
point is $g_0\exp(tY)g_0^{-1}g$.
Its scalar coefficient in the integrand is
\[
 F_m\bigl(P h'n'gP^{-1}\exp(tQ^{-1}YQ)\bigr).
 \tag{BSE.11}
\]
Here $Q^{-1}YQ$ is in the real Lie algebra of $G$, since
$Q^*H_{pk}Q=J_\epsilon$.
BST3 and BSP5 differentiate under the actual integral and identify
its derivative as $\ell(E(Q^{-1}YQ)m)$.

For the nilpotent directions define
\[
 \begin{split}
 \mathcal L(Y)&=\sum_{a=0}^{r-1}c_a\gamma_aY_{a,\operatorname{next}(a)},\\
 \gamma_a&=\begin{cases}1,&a+1<r,\\\sqrt{T_z},&a=r-1,\end{cases}
 \quad
 \operatorname{next}(a)=
       \begin{cases}a+1,&a+1<r,\\r+k,&a=r-1.\end{cases}
 \end{split}
 \tag{BSE.12}
\]
The Lie algebra of $N$ consists of the $H_{pk}$-skew-Hermitian
matrices strictly above the levels (BSE.7), by BWP1/BNV11.
Their consecutive-level exponential entries are their linear
entries, since every higher power raises level by at least two.
The defining frame fixes the outer coordinates and sends
$\sqrt{T_z}e_{r+k}$ to $z$.  Consequently the character of the
framed one-parameter subgroup is exactly
\[
 \psi\bigl(g_0\exp(tY)g_0^{-1}\bigr)
                         =\exp\bigl(it\,\operatorname{Re}\mathcal L(Y)\bigr).
 \tag{BSE.13}
\]
Differentiating BSE1 along this subgroup and comparing with
(BSE.11) gives
$\ell(E(Q^{-1}YQ)m)=i\operatorname{Re}\mathcal L(Y)\ell(m)$.
For the compact block, BSE1 similarly gives
$\ell(E(Q^{-1}\operatorname{emb}(A)Q)m)=d\tau(A)\ell(m)$
for skew-Hermitian $A$.  Complexification yields the second
identity of (BSE.8).

We make the nilpotent complexification and every scalar explicit.
Let $Z^\sharp=H_{pk}Z^*H_{pk}$.  Strictly increasing levels are
preserved by this involution.  Decompose
$Z=(Z-Z^\sharp)/2+i[-i(Z+Z^\sharp)/2]$ into two real
nilpotent Lie directions.  The complexified character is
\[
 \chi(Z)=\frac i2
       \left(\mathcal L(Z)-\overline{\mathcal L(Z^\sharp)}\right).
 \tag{BSE.14}
\]
Thus $\chi_{ij}=\chi(e_{ij})$ proves the first identity of (BSE.8).
The form $H_{pk}$ pairs $e_i$ with $e_{d-i}$ outside the middle
block, is $+1$ on $W$, and is $-1$ on $e_{r+k}$.
If $\iota$ fixes middle indices and interchanges the paired outer
indices, and $s_i=-1$ only at $i=r+k$, then
$e_{ij}^\sharp=s_is_j e_{\iota(j),\iota(i)}$.
Substituting in (BSE.12)--(BSE.14) shows that no edge outside
(BSE.9) contributes.  For $a<r-1$ the paired values are
$ic_a/2$ and $-i\overline{c_a}/2$.
For the last left edge $(r-1,r+k)$ and its right partner
$(r+k,r+k+1)$ they are
$i\sqrt{T_z}c_{r-1}/2$ and
$i\sqrt{T_z}\overline{c_{r-1}}/2$.
Their product is (BSE.10).
This also proves the source's nonzero \texttt{kappa} condition
with all complex $c_a\ne0$, without choosing their phases.
For $r=0$ there are no strictly increasing levels, so the nilpotent
conditions are empty.  BSE2 supplies (BSE.5) in all cases.
\end{proof}
% END BSE3

\subsection{Restriction and change of frame for the actual minors}

% BEGIN BSE4
\paragraph{BSE4. The whole Capelli--Whittaker restriction, including rank zero.}
\label{claim:BSE4}
Let $M,R$ be arbitrary complex modules, $E$ a matrix-unit action
on $M$ indexed by $0,\ldots,k+2r$, and
$l:M\to R$ a linear map.  Suppose
$lE_{ij}=\chi_{ij}l$ whenever $b(i)<b(j)$, and on these
strictly increasing-level pairs suppose $\chi_{ij}=0$ off the
edges (BSE.9).  No condition is imposed on the remaining values
of $\chi$.  Suppose also
$lE_{r+i,r+j}=D_{ij}l$ for $i,j<k$.
For $r>0$ and every $t\in\mathbb C$,
\[
 l\,C^E_{\widehat d,\widehat0}(t)
       =(-1)^k\kappa\,\mathcal C_D(t-r)\,l,
       \qquad \kappa=\prod_{\text{edges in (BSE.9)}}\chi_{ij}.
 \tag{BSE.15}
\]
This is the whole original \texttt{CapelliBStatement}; no nonzero
assumption on $\kappa$ is needed for the identity.
For $r=0$ one has, under the compact-block intertwining alone,
\[
                   l\,C^E_{\widehat k,\widehat k}(t)
                                =\mathcal C_D(t)\,l.
 \tag{BSE.16}
\]
The latter is the entire original \texttt{WhittakerZeroStatement}.

\begin{proof}
Use the exact shifted minors and expansions proved in BST6--BST8.
First take the minor with rows $0,\ldots,d-1$ and columns
$1,\ldots,d$, and expand at its last column, moving that column
to the front before the ordinary column expansion.
Every row has smaller level than the last column, so applying $l$
replaces each first operator by its character scalar.
All terms except the edge $(d-1,d)$ vanish.
Its row and column positions are both the last position, so its
sign is $+1$.  Removing them lowers the minor parameter by one.
Repeat for the $r$ successive right edges.  The result is their
product times the minor with rows $0,\ldots,r+k-1$, columns
$1,\ldots,r+k$, and parameter $t-r$.

Next expand at the first row using BST7's row expansion.
Until row $r-1$ is reached, only its next left edge survives.
It occupies the first column, so its sign is $+1$.
A first-row expansion leaves the remaining minor's parameter
unchanged.  After these $r-1$ removals the row list is
$r-1,r,\ldots,r+k-1$ and the column list is
$r,r+1,\ldots,r+k$, still at parameter $t-r$.
The only allowed edge from its first row is $(r-1,r+k)$,
which now occupies column position $k$.  Its sign is $(-1)^k$.
Removing it leaves the identical row and column lists
$r,\ldots,r+k-1$.

All discarded entries were off-diagonal, so their scalar shifts
were zero; hence the preceding substitutions did not alter any
Capelli parameter.  The collected scalar is exactly
$(-1)^k\kappa$.  On the remaining block,
$lE_{r+i,r+j}=D_{ij}l$ transfers every ordered factor through $l$,
and hence transfers its full determinant to
$\mathcal C_D(t-r)l$.  This proves (BSE.15), including $k=0$,
when the last determinant is the empty product.
For $r=0$ the minor in (BSE.16) is already the compact block;
the same factor-by-factor substitution proves it directly.
That argument does not even require bracket relations for the
rank-zero restriction identity.
\end{proof}
% END BSE4

% BEGIN BSE5
\paragraph{BSE5. Exact change of frame for every complementary minor.}
\label{claim:BSE5}
For an arbitrary associative complex algebra, elements $E_{ij}$
satisfying
$[E_{ab},E_{cd}]=\delta_{bc}E_{ad}-\delta_{da}E_{cb}$,
and an invertible complex matrix $Q$, define $E^Q$ by
(BSE.6).  Then for all $b,c$ and all $t$,
\[
 (-1)^{b+c}C^{E^Q}_{\widehat b,\widehat c}(t)
 =\sum_{a,j}Q_{ba}(Q^{-1})_{jc}(-1)^{a+j}
                                      C^E_{\widehat a,\widehat j}(t).
 \tag{BSE.17}
\]
In particular, with the literal $\texttt{nnPlus}$ and
$\texttt{nnMinus}$ of BST9,
\[
 \begin{split}
 \sum_a Q_{ba}\,
       \texttt{nnPlus}_E(Q^{-1}e_c,\omega,v)_a
      &=(-1)^{b+c}C^{E^Q}_{\widehat b,\widehat c}(\omega-1)v,\\
 \sum_j(Q^{-1})_{jc}\,
       \texttt{nnMinus}_E(e_b^tQ,\omega,v)_j
      &=(-1)^{b+c}C^{E^Q}_{\widehat b,\widehat c}(\omega)v.
 \end{split}
 \tag{BSE.18}
\]
The identities retain the source's ordered factors and complement
signs.
This is the whole original \texttt{CapelliCovStatement}:
its two square scalar matrices with product one are mutual
inverses, so take them to be $Q^{-1}$ and $Q$.

\begin{proof}
The entry relation, including its scalar shift, is
\[
 E^Q_{ij}+t\delta_{ij}
    =\sum_{a,b}(Q^{-1})_{ai}Q_{jb}
                                      (E_{ab}+t\delta_{ab}).
 \tag{BSE.19}
\]
The scalar part is correct because $QQ^{-1}=I$.
Expanding a shifted column minor in these entries and grouping
row choices gives
\[
 C^{E^Q}_{R,C}(t)
 =\sum_{\substack{I\text{ increasing}\\ |I|=p}}
   \ \sum_{j_0,\ldots,j_{p-1}}
       \det(Q^{-1}[I,R])
       \prod_{\ell=0}^{p-1}Q_{C_\ell,j_\ell}\,
                             C^E_{I,(j_0,\ldots,j_{p-1})}(t).
 \tag{BSE.20}
\]
Here is why this expansion is valid despite noncommutativity.
For a fixed ordered row choice, the original alternating
permutation sum gives the ordinary determinant of its scalar
$Q^{-1}$ entries.  Repeated row choices have determinant zero.
For distinct choices, group them by their increasing list $I$;
their permutation signs give exactly the row-alternating sum
defining $C^E_{I,(j_0,\ldots,j_{p-1})}$.
Only scalar factors were moved past operator factors.
Thus the operator products still have the original order and
the original descending shifts.

By the column alternation proved in BST6, repeated $j$'s now
vanish.  Grouping their permutations gives a second ordinary
determinant and yields the precise Cauchy--Binet formula
\[
 C^{E^Q}_{R,C}(t)
    =\sum_{\substack{I,J\text{ increasing}\\|I|=|J|=p}}
       \det(Q^{-1}[I,R])\,\det(Q[C,J])\,C^E_{I,J}(t).
 \tag{BSE.21}
\]
Apply this with $p=n-1$, $R=\widehat b$, $C=\widehat c$.
The ordinary cofactor identities give
\[
 \det(Q^{-1}[\widehat a,\widehat b])
      =(-1)^{a+b}\det(Q)^{-1}Q_{ba},\qquad
 \det(Q[\widehat c,\widehat j])
      =(-1)^{c+j}\det(Q)(Q^{-1})_{jc}.
 \tag{BSE.22}
\]
Substitute them into (BSE.21); the determinants cancel and give
(BSE.17).  Expanding the definitions of $\texttt{nnPlus}$ and
$\texttt{nnMinus}$ then gives (BSE.18), at parameters
$\omega-1$ and $\omega$ respectively.
No commutative adjugate identity was applied to the operator entries.
\end{proof}
% END BSE5

% BEGIN BSE6
\paragraph{BSE6. Both tensor periods in the same frame.}
\label{claim:BSE6}
At the fixed original period point $g$, put
\[
 j_*=\begin{cases}0,&r>0,\\d,&r=0,\end{cases}
 \quad
 b_+=\begin{cases}1,&r>0,\\-\sqrt{T_z},&r=0,\end{cases}
 \quad
 b_-=\begin{cases}1,&r>0,\\\sqrt{T_z},&r=0.\end{cases}
 \tag{BSE.23}
\]
For every tuple $f\in S^n$ the actual standard and dual periods are
\[
 \mathcal P_{Q_f}(g)=b_+\sum_jQ_{dj}\ell(f_j),\qquad
 \mathcal P_{Q_f^\vee}(g)=b_-\sum_j(Q^{-1})_{jj_*}\ell(f_j).
 \tag{BSE.24}
\]

\begin{proof}
Every element of $H_zN$ fixes $x_*$ from (BSE.2), by BSP4.
Because it is also $H_{pk}$-unitary, it fixes the row $x_*^*H_{pk}$.
Consequently the scalar multipliers in the standard and dual
tensor maps are constant along the left $H_zN$ factor of the
period.  Pulling them out and using the finite linear sum yields
\[
 \begin{split}
 \mathcal P_{Q_f}(g)
     &=\sum_j(x_*^*H_{pk}gP^{-1})_j\,\ell(f_j),\\
 \mathcal P_{Q_f^\vee}(g)
     &=\sum_j(Pg^{-1}x_*)_j\,\ell(f_j).
 \end{split}
 \tag{BSE.25}
\]
The individual summands are integrable by BST5 and BSP5, so this
linearity uses the actual integrals.

Now $gP^{-1}=g_0Q$.
If $r>0$, $x_*=g_0e_0$,
$e_0^*H_{pk}=e_d^t$, and $Pg^{-1}g_0=Q^{-1}$.
Thus the two multiplier lists are row $d$ of $Q$ and column
$0$ of $Q^{-1}$.
If $r=0$, $x_*=\sqrt{T_z}g_0e_d$ and
$e_d^*H_{pk}=-e_d^t$.  The lists are
$-\sqrt{T_z}$ times row $d$ of $Q$ and
$\sqrt{T_z}$ times column $d$ of $Q^{-1}$.
These are exactly (BSE.23)--(BSE.24).
\end{proof}
% END BSE6

\subsection{The two nonzero evaluated periods}

% BEGIN BSE7
\paragraph{BSE7. The whole original \texttt{EvalPStatement}, with an explicit witness.}
\label{claim:BSE7}
For every original datum fixed above, every $v\in V$, and every
index $s$, suppose $\ell(v)\ne0$ and
$u_j\ne\lambda_s+1/2$ for every $j<k$.
Then the exact original standard tensor period is nonzero with
\[
       y=Q^{-1}e_{j_*},\qquad
       \omega=-\lambda_s+\frac d2 .
 \tag{BSE.26}
\]
More precisely it equals
\[
 \mathcal P_{Q_{\texttt{nnPlus}(y,\omega,v)}}(g)
   =
 \begin{cases}
 \displaystyle \kappa\prod_{j<k}(u_j-\lambda_s-1/2)\,\ell(v),
                                                   &r>0,\\[2mm]
 \displaystyle -\sqrt{T_z}\prod_{j<k}(u_j-\lambda_s-1/2)\,\ell(v),
                                                   &r=0 .
 \end{cases}
 \tag{BSE.27}
\]
This proves the original statement for arbitrary original
$x_0,w,v,\xi,g$, not just for a special Bessel character or frame.

\begin{proof}
Every complementary-minor component of the tuple lies in $S$,
since it is a finite polynomial in its stable matrix-unit
operators applied to $v$.  Thus BSE6 applies.
Its first formula and the first identity of (BSE.18), with
$b=d,c=j_*$, give
\[
 \mathcal P_{Q_{\texttt{nnPlus}}}(g)
       =b_+(-1)^{d+j_*}
                 \ell\!\left(C^{E^Q}_{\widehat d,\widehat{j_*}}
                                           (\omega-1)v\right).
 \tag{BSE.28}
\]
The restriction of $E^Q$ to $S$ satisfies every hypothesis of
BSE4 by the actual functional law BSE3.
If $r>0$, $j_*=0$, so the right side of (BSE.28) is
\[
       (-1)^{d+k}\kappa\,p_u(\omega-1-r)\ell(v)
                              =\kappa\,p_u(\omega-1-r)\ell(v).
 \tag{BSE.29}
\]
The sign is one because $d+k=2r+2k$ is even.
If $r=0$, $j_*=d=k$, the complement is the compact principal
minor and $(-1)^{d+j_*}=1$.
BSE4 instead gives $-\sqrt{T_z}p_u(\omega-1)\ell(v)$.
In either case
\[
 \omega-1-r+u_j-\frac{k-1}{2}
                                  =u_j-\lambda_s-\frac12,
 \tag{BSE.30}
\]
where $r=0$ is included.  This proves the exact formula (BSE.27).
The scalar $\kappa$ is nonzero by (BSE.10), and $\sqrt{T_z}>0$.
Each displayed linear factor is nonzero by the original exclusion.
The empty product when $k=0$ is one.
Since $\ell(v)\ne0$, the period is nonzero.
All factors came from the prescribed integral; no abstract
branching functional was substituted.
\end{proof}
% END BSE7

% BEGIN BSE8
\paragraph{BSE8. The whole original \texttt{EvalMStatement}, with an explicit witness.}
\label{claim:BSE8}
For the same arbitrary original data, suppose $\ell(v)\ne0$ and
$u_j\ne\lambda_s-1/2$ for every $j<k$.
Take the row
\[
                    \alpha=e_d^tQ,\qquad
                    \omega=-\lambda_s+\frac d2 .
 \tag{BSE.31}
\]
Then the exact original dual tensor period is nonzero, and equals
\[
 \mathcal P_{Q^\vee_{\texttt{nnMinus}(\alpha,\omega,v)}}(g)
   =
 \begin{cases}
 \displaystyle \kappa\prod_{j<k}(u_j-\lambda_s+1/2)\,\ell(v),
                                                    &r>0,\\[2mm]
 \displaystyle \sqrt{T_z}\prod_{j<k}(u_j-\lambda_s+1/2)\,\ell(v),
                                                    &r=0 .
 \end{cases}
 \tag{BSE.32}
\]

\begin{proof}
As before the tuple lies in $S^n$.
Use the dual formula of BSE6 and the second identity of (BSE.18):
\[
 \mathcal P_{Q^\vee_{\texttt{nnMinus}}}(g)
       =b_-(-1)^{d+j_*}
               \ell\!\left(C^{E^Q}_{\widehat d,\widehat{j_*}}
                                                   (\omega)v\right).
 \tag{BSE.33}
\]
There is no minus-one shift in this complementary minor.
For $r>0$, BSE4 gives
$\kappa p_u(\omega-r)\ell(v)$ after the same sign cancellation.
For $r=0$ it gives
$\sqrt{T_z}p_u(\omega)\ell(v)$.
Now
\[
 \omega-r+u_j-\frac{k-1}{2}
                                  =u_j-\lambda_s+\frac12.
 \tag{BSE.34}
\]
This proves (BSE.32).  Its scalar is nonzero by the exact
negative-step exclusion, (BSE.10), and $T_z>0$.
The coefficient slot, character and inverse compact action are
unchanged from (BSE.1); this is precisely the original
\texttt{EvalMStatement}.
\end{proof}
% END BSE8

% BEGIN BSE9
\paragraph{BSE9. The exact remaining quotient input to the full Bessel theorem.}
\label{claim:BSE9}
Retain the complete original \texttt{QuotPStatement} and
\texttt{QuotMStatement}.  Then the reviewed DC/BNV/BSP results,
BST1--BST11 and BSE1--BSE8 imply the whole original
\texttt{StepStatement} and therefore the whole original
\texttt{BesselNVGap}.
On the inherited supplied \texttt{Forms0} scope, they also imply
the original \texttt{ConcreteSettingExists}.

\begin{proof}
BSE7 and BSE8 supply the two evaluation premises retained in BST11,
with all their original quantifiers, actual periods and exclusions.
Thus only its two quotient premises remain, and the assumed
\texttt{QuotP}/\texttt{QuotM} give both directions of
\texttt{StepStatement} by BST11.
BNV23 then supplies the arbitrary-data endomorphism-valued Bessel
nonvanishing theorem and the inherited BWP24 global-setting join.
The latter still applies on precisely the supplied
\texttt{Forms0} scope.
No new ambient form supplier or orbital/fiber comparison is used.
\end{proof}
% END BSE9

\paragraph{Remaining construction.}
The whole \texttt{QuotP} and \texttt{QuotM} statements are not
proved here.  They must still construct the literal spaces
$W_f,W'$, the exact image relation, the shifted lowest compact type
and Casimir, actual compact-group stability, and injectivity of
the period on the chosen quotient.
The now proved nonzero evaluations and scalar Casimir identities
supply their source premises; they do not by themselves construct
these spaces.  BSE9 is therefore an explicitly conditional
implication, and no whole selected global obligation is declared
closed by this chapter.
